\chapter{Gauduchon traces and the topology of potentials}
\label{chap:compactness}

Throughout this chapter, $\theta$ is a smooth real $(1,1)$-form on $X$, $\mathrm{d}V$ is a smooth positive probability volume form used for $L^1$ norms, and $\gamma$ is a Gauduchon metric, that is, a Hermitian form with $\ddc(\gamma^{n-1})=0$ (\cref{def:gauduchon-metric}). In dimension one every Hermitian form is Gauduchon, since $\gamma^{n-1}=1$.

The Monge--Amp\`ere measure of a $\theta$-plurisubharmonic function can have infinite mass, as we shall see in \cref{chap:products}, but its linear part is always under control. By \cref{prop:gauduchon-trace}, the positive measure $(\theta+\ddc u)\wedge\gamma^{n-1}$ has the same total mass for every $u\in\PSH(X,\theta)$. Since $f\mapsto\ddc f\wedge\gamma^{n-1}$ is a second-order elliptic operator whose kernel consists of the constants, this trace measure determines $u$ up to an additive constant. We show that it also determines the topology: on mean-normalized potentials, weak convergence of the trace measures is equivalent to $L^1$ convergence of the potentials. Testing the trace measures against Lipschitz functions turns this equivalence into an explicit complete metric, and the inverse of the elliptic operator compares this metric with the $L^1$ norm by sharp power bounds. Finally, the smooth form $\gamma^{n-1}$ can be replaced by a positive pluriharmonic current when the potential is continuous; the product then involves first derivatives of the potential, which are controlled by an energy estimate.

\Needspace{13\baselineskip}
\section{The trace map}
\label{sec:gauduchon-trace-topology}

Adding a constant leaves the trace unchanged. We therefore use the mean normalization
\[
 \mathcal K^0_{\theta,\mathrm{d}V}
   =\left\{u\in\PSH(X,\theta):\int_Xu\,\mathrm{d}V=0\right\}.
\]
For such a potential put
\[
 \mathcal T_\gamma(u)
     =(\theta+\ddc u)\wedge\gamma^{n-1},
 \qquad M_{\theta,\gamma}=\int_X\theta\wedge\gamma^{n-1}.
\]
By \cref{prop:gauduchon-trace}, the measure $\mathcal T_\gamma(u)$ is positive and has mass $M_{\theta,\gamma}$. In dimension one this follows from Stokes' theorem. For $f\in C^\infty(X,\mathbb R)$, set $Af=(\ddc f\wedge\gamma^{n-1})/\mathrm{d}V$.

\Needspace{13\baselineskip}
\begin{theorem}\label{thm:gaud-trace-homeomorphism}\index{Gauduchon trace!topology of potentials}
The map $\mathcal T_\gamma$ is an affine homeomorphism from the compact convex space $\mathcal K^0_{\theta,\mathrm{d}V}$, with its $L^1$ topology, onto a compact convex subset of the positive measures of mass $M_{\theta,\gamma}$, with their weak topology. For $u_j,u\in\mathcal K^0_{\theta,\mathrm{d}V}$,
\[
 u_j\longrightarrow u\text{ in }L^1
 \quad\Longleftrightarrow\quad
 \mathcal T_\gamma(u_j)\rightharpoonup\mathcal T_\gamma(u).
\]
If $M_{\theta,\gamma}=0$, both spaces consist of a single point.
\end{theorem}
\begin{proof}
The mean-normalized space is convex. It is the image of the compact space in \cref{thm:normalized-compactness} under the continuous surjection $v\mapsto v-\int_Xv\,\mathrm{d}V$, hence compact.

If $u_j\to u$ in $L^1$, their trace measures converge in distributions. Their common finite mass upgrades this to weak convergence, by uniform approximation of continuous tests by smooth functions.

For injectivity, equality of traces gives $A(u-v)=0$. Elliptic regularity for $A$ makes $u-v$ smooth, and the maximum principle makes it constant. Its mean is zero, so it vanishes. The potentials agree everywhere by \cref{lem:qpsh-almost-everywhere-order}.

A continuous injection from a compact space to the Hausdorff space of finite measures is a homeomorphism onto its image. Affineness gives convexity of the image. If its common mass is zero, every trace is the zero measure, and injectivity makes the potential space a singleton. In this case \cref{lem:zero-gauduchon-degree} also gives $\theta+\ddc u=0$ and smoothness of the unique normalized potential.
\end{proof}

Admissibility for the scalar inverse requires positivity of the full current. For a finite signed measure $\sigma$ of total mass zero, let $G_\gamma\sigma$ denote the mean-zero distribution solving
\begin{equation}\label{eq:gaud-trace-green}
 \ddc(G_\gamma\sigma)\wedge\gamma^{n-1}=\sigma.
\end{equation}
Indeed, $\ker A=\R$ by the maximum principle, and $A$ has index zero because its linear homotopy to a Laplace operator stays elliptic. Let $A^*$ be its formal adjoint with respect to $\mathrm{d}V$. Stokes' theorem gives $A^*1=0$, so $\ker A^*=\R$ as well. The Fredholm alternative solves the equation for smooth data exactly when $\int_X\sigma=0$. Duality with the continuous inverse of $A^*$ on smooth functions of mean zero extends $G_\gamma$ to top-degree distributions of total mass zero, including finite signed measures; imposing zero mean fixes the additive constant.

\begin{proposition}\label{prop:gaud-trace-admissible-data}
\index{Gauduchon trace!admissible measures}
A positive measure $\nu$ of mass $M_{\theta,\gamma}$ belongs to the image of $\mathcal T_\gamma$ if and only if the distribution
\[
 w=G_\gamma\bigl(\nu-\theta\wedge\gamma^{n-1}\bigr)
\]
satisfies $\theta+\ddc w\geq0$. In that case its canonical representative is the corresponding member of $\mathcal K^0_{\theta,\mathrm{d}V}$.
\end{proposition}
\begin{proof}
If $\nu=\mathcal T_\gamma(u)$, uniqueness in \eqref{eq:gaud-trace-green} gives $w=u$ as distributions. Conversely, \cref{prop:local-correction} gives a $\theta$-plurisubharmonic representative of $w$. Its mean is zero by construction, and \eqref{eq:gaud-trace-green} identifies its trace with $\nu$.
\end{proof}

In dimension one the positivity condition is automatic: $\theta+\ddc w=\nu$. Thus every positive measure of the specified mass occurs. In higher dimension, the scalar trace condition must be supplemented by positivity of the full current.

\section{A metric on normalized potentials}
\label{sec:gauduchon-trace-metric}

Testing trace differences against bounded Lipschitz functions gives a complete metric, which the elliptic inverse compares quantitatively with the $L^1$ norm. Fix a distance $d$ induced by a smooth Riemannian metric on $X$. For a real Lipschitz function $f$ on $X$, write $\operatorname{Lip}_d(f)=\sup_{x\ne y}|f(x)-f(y)|/d(x,y)$.

\begin{corollary}\label{cor:gaud-trace-metric}\index{Gauduchon trace!metric}
For $u,v\in\mathcal K^0_{\theta,\mathrm{d}V}$, the formula
\[
 d_{\gamma,d}(u,v)
 =\sup_{\substack{\|f\|_\infty\leq1\\\operatorname{Lip}_d(f)\leq1}}
       \left|\int_X f\,\mathrm{d}\bigl(\mathcal T_\gamma(u)-\mathcal T_\gamma(v)\bigr)\right|
\]
is a bounded complete metric on $\mathcal K^0_{\theta,\mathrm{d}V}$, inducing its $L^1$ topology. Changing $\gamma$ or $d$ gives the same topology and a uniformly equivalent metric on this compact space.
\end{corollary}
\begin{proof}
Lipschitz functions are uniformly dense in continuous functions, and the test family in the supremum is uniformly bounded and equicontinuous. The first fact separates measures and gives weak convergence from metric convergence. The Arzel\`a--Ascoli theorem supplies finite uniform nets for the test family and gives the converse. Symmetry and the triangle inequality follow directly from the formula, and its value is at most $2M_{\theta,\gamma}$.

Pull back this metric by the homeomorphism in \cref{thm:gaud-trace-homeomorphism}. The resulting metric space is compact, hence complete. For two choices of the auxiliary data, both metrics induce the same compact $L^1$ topology, so the identity maps in both directions are uniformly continuous.
\end{proof}

\begin{proposition}\label{prop:gaud-trace-quantitative}
\index{Gauduchon trace!quantitative bounds}
There is a constant $C\geq1$, depending only on $X,\theta,\gamma,\mathrm{d}V$ and $d$, such that
\[
 C^{-1}\|u-v\|_{L^1}
 \leq d_{\gamma,d}(u,v)
 \leq C\|u-v\|_{L^1}^{1/2}
 \qquad(u,v\in\mathcal K^0_{\theta,\mathrm{d}V}).
\]
\end{proposition}
\begin{proof}
Put $w=u-v$, $\sigma=\mathcal T_\gamma(u)-\mathcal T_\gamma(v)$ and $M=M_{\theta,\gamma}$, and write $\|\sigma\|_{\mathrm{TV}}$ for the total variation of $\sigma$. Then $\int_Xw\,\mathrm{d}V=0$, $\sigma=\ddc w\wedge\gamma^{n-1}$ and $\|\sigma\|_{\mathrm{TV}}\leq2M$. If $M=0$, \cref{thm:gaud-trace-homeomorphism} gives $u=v$, so both estimates hold.

Since $\ker A=\ker A^*=\R$, for a smooth real function $g$, the equation
\[
 A^*\phi=g-\int_Xg\,\mathrm{d}V,\qquad
 \int_X\phi\,\mathrm{d}V=0
\]
has a unique smooth solution. For $p>2n$, elliptic $W^{2,p}$ estimates on mean-zero functions and Sobolev embedding give $\|\phi\|_\infty+\operatorname{Lip}_d(\phi)\leq C\|g\|_\infty$.
\Needspace{6\baselineskip}
Pairing the distributional equation $Aw\,\mathrm{d}V=\sigma$ with $\phi$ yields
\[
 \left|\int_Xwg\,\mathrm{d}V\right|
 =\left|\int_X\phi\,\mathrm{d}\sigma\right|
 \leq C\|g\|_\infty d_{\gamma,d}(u,v).
\]
Taking the supremum over smooth real $g$ with $\|g\|_\infty\leq1$ gives $\|w\|_{L^1}$ and proves the first estimate.

For the second estimate, let $f$ be any test function in the definition of $d_{\gamma,d}$. A fixed finite coordinate cover, smooth partition of unity and convolution at scale $0<\epsilon<\epsilon_0$ give smooth functions $f_\epsilon$ satisfying
\[
 \|f-f_\epsilon\|_\infty\leq C\epsilon,
 \qquad \|A^*f_\epsilon\|_\infty\leq C\epsilon^{-1}.
\]
The constants are uniform over this test family. To see the second bound, localize $f$ by the fixed partition. The localized functions have uniformly bounded values and first weak derivatives. In each second derivative of their convolution, move one derivative onto the mollifier; its $L^1$ norm is $O(\epsilon^{-1})$. The lower-order terms of $A^*$ have fixed smooth coefficients and satisfy the same bound. Hence
\[
 \left|\int_Xf\,\mathrm{d}\sigma\right|
 \leq C\bigl(M\epsilon+\epsilon^{-1}\|w\|_{L^1}\bigr).
\]
If $0<\|w\|_{L^1}<\epsilon_0^2$, choose $\epsilon=\|w\|_{L^1}^{1/2}$. For larger norms use $d_{\gamma,d}(u,v)\leq2M$, and for $w=0$ use injectivity. Taking the supremum over $f$ proves the asserted bound.
\end{proof}

\begin{example}\label{ex:gaud-trace-sharp-exponents}
\index{Gauduchon trace!sharp exponents}
Let $X=\mathbb C/(\mathbb Z+\mathrm{i}\mathbb Z)$, with $z=x+\mathrm{i}y$, $\theta=\gamma=\eta=(\mathrm{i}/2)\mathrm{d}z\wedge\mathrm{d}\bar z$, $\mathrm{d}V=\eta$ and the flat distance $d$. For integers $k\geq1$, set
\[
 u_k=\frac{\cos(2\pi kx)}{2\pi k^2}.
\]
These smooth mean-zero potentials satisfy
\[
 \theta+\ddc u_k=\left(1-\frac12\cos(2\pi kx)\right)\eta>0,
 \qquad \|u_k\|_{L^1}=\frac1{\pi^2k^2}.
\]
The function $f_k=\cos(2\pi kx)/(2\pi k)$ has supremum norm at most one and Lipschitz constant at most one. It follows that
\[
 d_{\gamma,d}(u_k,0)
 \geq\left|\int_Xf_k\,\ddc u_k\right|
 =\frac1{8\pi k}.
\]
Consequently, the exponent $1/2$ in \cref{prop:gaud-trace-quantitative} cannot be replaced by any larger exponent, even for smooth strictly positive potentials on a fixed curve. The power one in the lower bound cannot be replaced by a smaller positive power: for $0<t\leq1$, both $\|tu_1\|_{L^1}$ and $d_{\gamma,d}(tu_1,0)$ equal $t$ times their positive values at $t=1$.
\end{example}

\section{Continuous potentials and pluriharmonic traces}
\label{sec:continuous-pluriharmonic-traces}

For $n\geq2$, let $\omega$ be a Hermitian form on $X$, and let $G$ be a positive $(n-1,n-1)$-current with $\ddc G=0$. Such currents include $\gamma^{n-1}$, but their coefficients may be singular measures. The equation $\ddc G=0$ converts uniform control of a potential difference into a gradient-energy bound. This bound permits a canonical trace measure for continuous potentials.

\begin{lemma}\label{lem:pluriharmonic-gradient-energy}\index{Pluriharmonic current!gradient energy}
Let $(X,\omega)$ be a compact Hermitian manifold of dimension $n\geq2$, and let $G$ be a positive $(n-1,n-1)$-current with $\ddc G=0$. Put $\mu=\omega\wedge G$ and $M=\mu(X)$. If $a,b$ are smooth real functions with $\ddc a,\ddc b\geq-C\omega$, where $C\geq0$, then
\[
 0\leq\int_X\mathrm{d}(a-b)\wedge\mathrm{d}^c(a-b)\wedge G
 \leq2C M\|a-b\|_\infty.
\]
For every smooth real function $v$, the total variation norms of the order-zero currents satisfy
\[
 \|\mathrm{d}v\wedge G\|_{\mathrm{TV}}+
 \|\mathrm{d}^cv\wedge G\|_{\mathrm{TV}}
 \leq C_\omega M^{1/2}
       \left(\int_X\mathrm{d}v\wedge\mathrm{d}^cv\wedge G\right)^{1/2}.
\]
Here the norms use the metric $\omega$, and $C_\omega$ depends only on that fixed geometry and the operator normalization.
\end{lemma}
\begin{proof}
Set $v=a-b$. The identity $\ddc(v^2)=2\mathrm{d}v\wedge\mathrm{d}^cv+2v\ddc v$ and $\ddc G=0$ give
\[
 \int_X\mathrm{d}v\wedge\mathrm{d}^cv\wedge G
 =-\int_Xv\ddc v\wedge G.
\]
The measures $(\ddc a+C\omega)\wedge G$ and $(\ddc b+C\omega)\wedge G$ are positive, each of mass $CM$, since their potential terms integrate to zero. Their difference is $\ddc v\wedge G$ and has total variation at most $2CM$. This proves the first estimate; nonnegativity follows from positivity of $\mathrm{d}v\wedge\mathrm{d}^cv=(\mathrm{i}/2\pi)\partial v\wedge\bar\partial v$.

In a unitary frame, positivity represents $G$ by a positive semidefinite matrix of measures. Its entries are absolutely continuous with respect to its trace measure $\mu$: the diagonal entries are dominated by the trace, and the off-diagonal entries satisfy the matrix Cauchy--Schwarz inequality. The resulting measurable matrix has trace bounded by a fixed normalization constant. Pairing $\mathrm{d}v\wedge G$ or $\mathrm{d}^cv\wedge G$ with a test one-form $\eta$ and applying the matrix inequality and then integral Cauchy--Schwarz gives
\[
 |\langle\mathrm{d}v\wedge G,\eta\rangle|
 +|\langle\mathrm{d}^cv\wedge G,\eta\rangle|
 \leq C_\omega\|\eta\|_\infty M^{1/2}
 \left(\int_X\mathrm{d}v\wedge\mathrm{d}^cv\wedge G\right)^{1/2}.
\]
A finite frame cover and partition of unity give the global estimate. Taking the supremum over smooth tests of norm at most one yields the total variation bound for these measure-valued currents. When $M=0$, positivity gives $G=0$, so both estimates hold.
\end{proof}

\begin{theorem}\label{thm:continuous-pluriharmonic-trace-product}\index{Pluriharmonic current!continuous-potential product}\index{Gauduchon trace!pluriharmonic current tests}
Let $(X,\omega)$ be a compact Hermitian manifold of dimension $n\geq2$, let $\alpha$ be a smooth real $(1,1)$-form, and let $u\in\PSH(X,\alpha)\cap C^0(X)$. For every positive $(n-1,n-1)$-current $G$ with $\ddc G=0$, the product
\[
 (\alpha+\ddc u)\wedge G
\]
is a canonically defined positive measure of mass $\int_X\alpha\wedge G$. It agrees with smooth multiplication if $u$ or $G$ is smooth. If $u_j,u\in\PSH(X,\alpha)\cap C^0(X)$ converge uniformly and positive pluriharmonic currents $G_j\rightharpoonup G$ converge weakly, then
\[
 (\alpha+\ddc u_j)\wedge G_j
 \rightharpoonup(\alpha+\ddc u)\wedge G.
\]
\end{theorem}
\begin{proof}
Choose smooth approximants $a_k\to u$ from \cref{lem:continuous-potential-smooth-loss}. There is a common $C>0$ such that $\ddc a_k\geq-C\omega$. Write $M=\int_X\omega\wedge G$. By \cref{lem:pluriharmonic-gradient-energy}, $\mathrm{d}a_k\wedge G$ and $\mathrm{d}^ca_k\wedge G$ are Cauchy in total variation. Denote their order-zero limits by $\mathrm{d}u\wedge G$ and $\mathrm{d}^cu\wedge G$. Comparing two uniformly convergent smooth sequences with a common lower Hessian bound proves independence of the sequence. Continuous multiplication also gives $a_kG\to uG$ in total variation.

Define the product by the distributional formula
\begin{equation}\label{eq:continuous-pluriharmonic-product}
 (\alpha+\ddc u)\wedge G
 =\alpha\wedge G+\mathrm{d}(\mathrm{d}^cu\wedge G)
   -\mathrm{d}^c(\mathrm{d}u\wedge G)-\ddc(uG).
\end{equation}
For smooth $u$, the full Leibniz identity gives this formula because $\ddc G=0$. Consequently $(\alpha+\ddc a_k)\wedge G$ converges to the displayed distribution. The measures $(\alpha+\ddc a_k+\varepsilon_k\omega)\wedge G$ are positive, with masses $\int_X\alpha\wedge G+\varepsilon_kM$. Their distributional limit is therefore a positive measure, and testing on the constant function one gives its asserted mass. The uniform mass bound upgrades distributional convergence to weak convergence of measures. If $G$ is smooth, uniform convergence of $a_k$ also identifies the limit by distributional multiplication with that smooth form.

Adding a smooth function $h$ to $u$ adds $\mathrm{d}h\wedge G$ and $\mathrm{d}^ch\wedge G$ to its gradient limits. Replacing $\alpha$ by $\alpha-\ddc h$ therefore leaves the product unchanged. If two continuous representations define the same $(1,1)$-current, their potential difference has smooth $\ddc$; taking its trace in a coordinate chart and applying elliptic regularity makes that difference smooth. Thus the product depends only on the current $\alpha+\ddc u$ and on $G$.

It remains to prove joint continuity. For each of $u$ and the $u_j$, retain a smooth approximating tail with curvature errors at most one. By \cref{lem:continuous-potential-smooth-loss} and a fixed bound $\alpha\leq(C-1)\omega$, all these approximants then satisfy the same lower Hessian bound $-C\omega$. The numbers $M_j=\int_X\omega\wedge G_j$ are bounded by weak convergence. Fix a smooth approximant $a_k$ of $u$. Apply the energy and total variation estimates to smooth approximants of $u_j$ and to $a_k$, then pass to their gradient limits. With $\delta_{j,k}=\|u_j-a_k\|_\infty$, the gradient differences have total variation at most $C' M_j\sqrt{\delta_{j,k}}$, where $C'$ is independent of $j,k$.

For a smooth real test function $\phi$, \eqref{eq:continuous-pluriharmonic-product} now gives
\[
 \left|\left\langle(\alpha+\ddc u_j)\wedge G_j
       -(\alpha+\ddc a_k)\wedge G_j,\phi\right\rangle\right|
 \leq C_\phi M_j(\sqrt{\delta_{j,k}}+\delta_{j,k}),
\]
where $C_\phi$ depends on the first two derivatives of $\phi$, the fixed geometry and $C$. The same estimate holds for $u,G$. For fixed $k$, smooth multiplication gives $(\alpha+\ddc a_k)\wedge G_j\rightharpoonup(\alpha+\ddc a_k)\wedge G$. First let $j\to\infty$, and then $k\to\infty$; since $u_j\to u$ uniformly, the error tends to zero. This proves distributional convergence. Positivity and the bounded masses $\int_X\alpha\wedge G_j$ upgrade it to weak convergence of measures.
\end{proof}

For $G=\gamma^{n-1}$, the product in \cref{thm:continuous-pluriharmonic-trace-product} is the Gauduchon trace of a continuous potential. Its gradient construction also allows singular pluriharmonic test currents.

